\documentclass[10pt,a4paper]{amsart}
\usepackage[margin=19mm]{geometry}
\usepackage{amsmath,amssymb,mathtools}
\usepackage[scr=rsfs,cal=euler]{mathalfa}
\usepackage{xfrac}
\usepackage[unicode,hidelinks,bookmarks=false]{hyperref}
\newcommand{\ie}{i.e.\ }
\newcommand{\cf}{cf.\ }
\newcommand{\resp}{respectively\ }
\newcommand{\Subsection}{\subsection}
\providecommand{\nobreakdash}{\nobreak\hbox{-}}
\setlength{\emergencystretch}{2em}
\multlinegap0pt
\usepackage{amssymb}
\newtheorem{theorem}{Theorem}
\title{Necessary and Sufficient Conditions for Existence of a Unique Solution to Gamma Moment Closure for the Stochastic Ricker Equation}
\author{Haiyan Wang, Melinda Wang}
\date{Source: arXiv:2601.20523v1 (CC BY 4.0)}
\begin{document}
\maketitle

\noindent\emph{Source and licence.} Necessary and Sufficient Conditions for Existence of a Unique Solution to Gamma Moment Closure for the Stochastic Ricker Equation. Version arXiv:2601.20523v1 (CC BY 4.0), distributed by arXiv under the Creative Commons Attribution 4.0 International licence, http://creativecommons.org/licenses/by/4.0/, as recorded in the arXiv licence metadata for this exact version and shown on the abstract page https://arxiv.org/abs/2601.20523v1. Full licence notice: \texttt{licenses/arxiv-2601.20523-CC-BY-4.0.txt}.\par\smallskip

\noindent\textit{Editorial scope.} Counted unit: the article's theorem auxi (source0.tex lines 356--368). The article's complete original proof of it is delivered with it (source0.tex lines 370--433, one \texttt{proof} environment of 64 lines). No mathematics is written, repaired or re-derived: the statement and the proof are the article's own lines, in the article's own notation, and no retained line is substituted or re-flowed.  Retained uncounted as the source-local context the proof needs: lines 83--85; lines 134--145; lines 153--170.  Nothing was omitted from any retained span, so the package carries no omission record.  No retained line contains a citation, so the wrapper carries no bibliography and no bibliography entry.  No retained line contains a figure environment, an image-inclusion command or drawing code, so no image is delivered and the package declares no asset dependency.  Editorial formatting only, in addition: an AMS \texttt{amsart} preamble that restates the article's own declarations and macros used by the retained lines, namely the environments \texttt{theorem} on separate counters, together with the standard packages those lines need.  The retained environments print in this document as Theorem 1 in order of appearance, whereas the article prints the same items with the numbering of its own full document.  The complete native source of the article as distributed in the arXiv e-print for this exact version is bundled unchanged as \texttt{sources/193526/source0.tex}.
\begin{equation}
  X_{n+1} = X_n e^{r(1 - X_n)} \varepsilon_n,\, n=1,2,... \label{eq:model}
\end{equation}

To obtain a tractable system for the stochastic difference equations, we adopt the Gamma moment-closure approximation. This approach does not assume that $X_n$ itself follows the Gamma distribution,  but rather that its first few moments satisfy the same relationships as those of a Gamma random variable. 
Specifically, we assume that the moments of $X_n, n=1,2,...$, satisfy
\begin{align}
\begin{aligned}
\mu_n \qquad\qquad & = \mathbb{E}[X_n]\\
 s_n \qquad\qquad &= \mathrm{Var}(X_n) \\
\mathbb{E}\big[X_n e^{-\tau X_n}\big] &= \frac{\mu_n}{\big(1 + \tau \theta_n\big)^{k_n+1}}, \\
\mathbb{E}\big[X_n^2 e^{-\tau X_n}\big] &= \frac{\mu_n^2+s_n}{(1+\tau \theta_n)^{k_n+2}}.
\end{aligned}
\label{r1}
\end{align}
where $ k_n=\frac{\mu_n^2}{s_n}, \, \theta_n=\frac{s_n}{\mu_n}$ and $\mu_n, s_n>0, \tau>0$.

\begin{theorem}[Gamma moment closure equations for the stochastic Ricker map] \label{differeceEq}
Consider the stochastic Ricker difference equation \eqref{eq:model}. Suppose that $\varepsilon_n$ 
is any random variable and independent of $X_n$ with $$E[\varepsilon_n] = 1, \,\,E[\varepsilon_n^2] = v > 1.$$ Assume that $X_n$ satisfies the Gamma moment-closure approximation \eqref{r1}.  Then its mean and variance \((\mu_n,s_n)\) updates in \eqref{eq:model}  can be written as follows:

\begin{equation}\label{eq:closed-mu-s}
(\mu_{n+1},s_{n+1}) = (G_1(\mu_n,s_n), G_2(\mu_n,s_n))
\end{equation}
where
\begin{align}
\begin{aligned}
\mu_{n+1} &= G_1(\mu_n,s_n)= e^{r}\,\frac{\mu_n}{\big(1 + r \theta_n\big)^{k_n+1}}, \\
s_{n+1} & =G_2(\mu_n,s_n)= v\,e^{2r}\,\frac{\mu_n^2+s_n}{(1+2r \theta_n)^{k_n+2}}
- e^{2r}\,\frac{\mu_n^2}{(1+r \theta_n)^{2(k_n+1)}}.
\end{aligned}
\label{expressions}
\end{align}
where $ k_n=\frac{\mu_n^2}{s_n}, \, \theta_n=\frac{s_n}{\mu_n}$.
\end{theorem}

\begin{theorem}[ Auxiliary function $F$]\label{auxi} Under the assumptions of Theorem \ref{differeceEq}. Assume that for $r>0, v>1$, and  define the following function of $z$ for \(z>0\): 
\begin{equation}\label{F1}
F(z;r,v)\;:=\;\Big(\frac{r}{\ln(1+z)}+1\Big)\ln(1+2z)\;-\;2r-\ln v,
\qquad z>0. 
\end{equation}
Then \((\mu,s)\) with $\mu>0,s>0$ is a feasible equilibrium  of Gamma moment-closure system \eqref{eq:closed-mu-s} if and only if
\begin{equation}\label{F2}
\mu \;=\; \frac{z\big(r-\ln(1+z)\big)}{r\ln(1+z)},
\qquad
s \;=\; \frac{\mu z}{r}.
\end{equation}
where $z>0$ is the solution of $F(z;r,v)=0$.
\end{theorem}

\begin{proof}
Assume $(\mu, s)$ with \(\mu>0,s>0\) is an equilibrium of \eqref{eq:closed-mu-s}. Note that $k=\frac{\mu^2}{s}, \theta=\frac{s}{\mu}$. From the first equilibrium condition of \eqref{eq:closed-mu-s} 
\[
\mu = e^{r}\,\frac{\mu}{(1 + r\theta)^{k+1}}
\]
we may cancel \(\mu\) and obtain
\[
(1 + r\theta)^{k+1} = e^{r}.
\]
Taking logarithms yields the identity
\begin{equation}\label{ad}
(k+1)\ln(1+r\theta) = r. 
\end{equation}
Introduce the auxiliary variable
\[
z := r\theta = \frac{r s}{\mu},
\]
so that \(1+r\theta=1+z\) and \(1+2r\theta=1+2z\).  Using \(k=\dfrac{\mu^2}{s}\) and \(\theta=\dfrac{s}{\mu}\) one finds
\[
k+1 = \mu \frac{\mu}{s}+1 = \mu \frac{r}{z}+1= \frac{\mu r}{z}+1.
\]
Thus equation \eqref{ad} becomes
\[
\Big(\frac{\mu r}{z}+1\Big)\ln(1+z) = r.
\]
Solving this for \(\mu r/z\) gives
\[
\frac{\mu r}{z} = \frac{r}{\ln(1+z)} - 1.
\]
Hence
\[
\mu \;=\; \frac{z}{r}\Big(\frac{r}{\ln(1+z)} - 1\Big)
= \frac{z\big(r-\ln(1+z)\big)}{r\ln(1+z)},
\]
which is the first relation in \eqref{F2}. Then \(s=\mu\theta=\mu z/r\), producing the second relation in  \eqref{F2}.

Next we derive \eqref{F1}.  From the first equilibrium condition of \eqref{eq:closed-mu-s}, we have
\[
(1 + r\theta)^{2(k+1)} = e^{2r},
\]
so the last term in the variance update simplifies at equilibrium:
\[
e^{2r}\,\frac{\mu^2}{(1 + r\theta)^{2(k+1)}} = \mu^2.
\]
Thus the variance equilibrium relation reduces to
\[
s = v e^{2r}\,\frac{\mu^2+s}{(1+2r\theta)^{k+2}} - \mu^2,
\]
or equivalently
\[
(\mu^2+s)\Big[1 - \frac{v e^{2r}}{(1+2r\theta)^{k+2}}\Big] = 0.
\]
Since \(\mu^2+s>0\) for any nontrivial equilibrium, we obtain
\begin{equation}\label{afd}
(1+2r\theta)^{k+2} = v e^{2r}.
\end{equation}
Since  $$1+2r\theta=1+2z$$  and by \eqref{ad}, it follows  
$$k+2 = (k+1)+1 = \dfrac{r}{\ln(1+z)}+1$$
Thus, taking logarithms and substituting in \eqref{afd}, we have 
\[
\Big(\frac{r}{\ln(1+z)}+1\Big)\ln(1+2z) = 2r + \ln v,
\]
which rearranges to the scalar equation \(F(z;r,v)=0\) displayed in \eqref{F1}. Now note that the derivation works for both the ``if" direction and the ``only if" direction. This completes the reduction: any positive equilibrium of  \eqref{eq:closed-mu-s} corresponds to a positive root \(z>0\) of \(F\) in  \eqref{F1}, and \((\mu,s)\) is then given by \eqref{F2}.
\end{proof}

\end{document}
